\frfilename{mt234.tex}
\versiondate{11.4.09}
\copyrightdate{2008}

\def\chaptername{Teorema Radon-Nikod\'ym}
\def\sectionname{Operasi pada ukuran}

\newsection{234}

Saya menggunakan beberapa halaman untuk menguraikan sejumlah konstruksi baku.
Gagasan-gagasannya langsung, tetapi sejumlah perincian perlu dikerjakan
agar dapat dipadukan secara kokoh ke dalam kerangka umum yang saya gunakan.
Langkah pertama adalah memperkenalkan secara formal fungsi-fungsi \imp\ 
(234A-234B), kelas transformasi terpenting di antara ruang-ruang ukur.
Untuk membangun ukuran baru, kita memiliki konsep
ukuran citra (234C-234E), %234C 234D 234E
jumlah ukuran (234G-234H), dan ukuran integral tak tentu
(234I-234O).  %234I 234J 234K 234L 234M 234N 234O
Terakhir, saya menyebutkan suatu cara mengurutkan ukuran-ukuran pada suatu
himpunan tertentu (234P-234Q).


\leader{234A}{Fungsi pelestari ukuran melalui
prapeta}\dvAformerly{2{}35G}\cmmnt{ Sudah tiba waktunya saya memperkenalkan
hal yang paling mendekati suatu `morfisme' dalam teori ukuran.}
Jika $(X,\Sigma,\mu)$ dan $(Y,\Tau,\nu)$ adalah ruang ukur, suatu fungsi
$\phi:X\to Y$ disebut {\bf \imp} jika
$\phi^{-1}[F]\in\Sigma$ dan $\mu(\phi^{-1}[F])=\nu F$ untuk setiap
$F\in\Tau$.

\leader{234B}{Proposisi} Misalkan $(X,\Sigma,\mu)$ dan $(Y,\Tau,\nu)$
ruang ukur, dan $\phi:X\to Y$ suatu fungsi \imp.


(a)\dvAformerly{2{}35Hc} Jika $\hat\mu$, $\hat\nu$ masing-masing adalah
pelengkapan dari $\mu$, $\nu$, maka $\phi$ juga \imp\ untuk $\hat\mu$ dan
$\hat\nu$.

(b)\dvAformerly{2{}35Xe} $\mu$ adalah ukuran peluang jika dan hanya jika
$\nu$ adalah ukuran peluang.

(c)\dvAformerly{2{}35Xe}
 $\mu$ berhingga total jika dan hanya jika $\nu$ berhingga total.

(d)\dvAformerly{2{}35Xe}(i)
Jika $\nu$ $\sigma$-hingga, maka $\mu$ $\sigma$-hingga.

\quad(ii) Jika $\nu$ semihingga dan $\mu$ $\sigma$-hingga, maka $\nu$
$\sigma$-hingga.

(e)\dvAformerly{2{}35Xe}(i)
Jika $\nu$ $\sigma$-hingga dan tanpa atom, maka $\mu$ tanpa atom.

\quad(ii) Jika $\nu$ semihingga dan $\mu$ atomik murni, maka
$\nu$ atomik murni.

(f)\dvAnew{2008}(i) $\mu^*\phi^{-1}[B]\le\nu^*B$ untuk
setiap $B\subseteq Y$.

\quad(ii) $\mu^*A\le\nu^*\phi[A]$ untuk setiap $A\subseteq X$.

(g)\dvAformerly{2{}35Hc}
 Jika $(Z,\Lambda,\lambda)$ adalah ruang ukur lain, dan
$\psi:Y\to Z$ bersifat \imp, maka $\psi\phi:X\to Z$ bersifat \imp.

\proof{{\bf (a)} Jika $\hat\nu$ mengukur $F$, terdapat $F'$, $F''\in\Tau$
sedemikian sehingga $F'\subseteq F\subseteq F''$ dan
$\nu(F''\setminus F')=0$. Sekarang

\Centerline{$\phi^{-1}[F']\subseteq\phi^{-1}[F]\subseteq\phi^{-1}[F'']$,
\quad$\mu(\phi^{-1}[F'']\setminus\phi^{-1}[F'])=\nu(F''\setminus F')=0$,}

\noindent sehingga $\hat\mu$ mengukur $\phi^{-1}[F]$ dan

\Centerline{$\hat\mu(\phi^{-1}[F])=\mu\phi^{-1}[F']=\nu F'=\hat\nu F$.}

\noindent Karena $F$ sembarang, $\phi$ bersifat \imp\ untuk $\hat\mu$ dan
$\hat\nu$.

\medskip

{\bf (b)-(c)} tentu saja jelas.

\medskip

{\bf (d)(i)} Jika $\sequencen{F_n}$ adalah penutup $Y$ oleh
himpunan-himpunan berukuran berhingga terhadap $\nu$, maka
$\sequencen{\phi^{-1}[F_n]}$ adalah penutup $X$ oleh himpunan-himpunan
berukuran berhingga terhadap $\mu$.

\medskip

\quad{\bf (ii)} Misalkan $\Cal F\subseteq\Tau$ suatu keluarga saling lepas
dari himpunan-himpunan yang tidak terabaikan terhadap $\nu$. Maka
$\family{F}{\Cal F}{\phi^{-1}[F]}$ adalah keluarga saling lepas dari
himpunan-himpunan yang tidak terabaikan terhadap $\mu$. Menurut 215B(iii),
$\Cal F$ terhitung. Dengan menggunakan 215B(iii) dalam arah sebaliknya,
$\nu$ bersifat $\sigma$-hingga.

\medskip

{\bf (e)(i)} Andaikan $E\in\Sigma$ dan $\mu E>0$. Misalkan
$\sequencen{F_n}$ suatu penutup $Y$ oleh himpunan-himpunan berukuran
berhingga terhadap $\nu$. Karena $\nu$ tanpa atom, untuk setiap $n$ kita
dapat menemukan suatu partisi berhingga $\family{i}{I_n}{F_{ni}}$
dari $F_n$ sedemikian sehingga $\nu F_{ni}<\mu E$ untuk setiap $i\in I_n$
(gunakan 215D berulang kali). Sekarang
$X=\bigcup_{n\in\Bbb N,i\in I_n}\phi^{-1}[F_{ni}]$, sehingga terdapat
$n\in\Bbb N$ dan $i\in I_n$ dengan

\Centerline{$0<\mu(E\cap\phi^{-1}[F_{ni}])
\le\mu\phi^{-1}[F_{ni}]=\nu F_{ni}<\mu E$,}

\noindent dan $E$ bukan atom terhadap $\mu$. Karena $E$ sembarang, $\mu$
tanpa atom.

\medskip

\quad{\bf (ii)} Andaikan $F\in\Tau$ dan $\nu F>0$. Karena $\nu$
semihingga, terdapat suatu $F_1\subseteq F$ sedemikian sehingga
$0<\nu F_1<\infty$. Sekarang $\mu\phi^{-1}[F_1]>0$; karena $\mu$
atomik murni, terdapat suatu atom terhadap $\mu$,
$E\subseteq\phi^{-1}[F_1]$.

Misalkan $\Cal G$ adalah himpunan semua $G\in\Tau$ sedemikian sehingga
$G\subseteq F_1$ dan $\mu(E\cap\phi^{-1}[G])=0$. Gabungan setiap barisan
di dalam $\Cal G$ termasuk dalam $\Cal G$, sehingga menurut 215Ac terdapat
suatu $H\in\Cal G$ sedemikian sehingga $\nu(G\setminus H)=0$ kapan pun
$G\in\Cal G$. Tinjau $F_1\setminus H$. Kita mempunyai

\Centerline{$\nu(F_1\setminus H)
=\mu(\phi^{-1}[F_1]\setminus\phi^{-1}[H])
\ge\mu(E\setminus\phi^{-1}[H])
=\mu E>0$.}

\noindent Jika $G\in\Tau$ dan $G\subseteq F_1\setminus H$,
maka salah satu dari $E\cap\phi^{-1}[G]$ dan
$E\setminus\phi^{-1}[G]$ terabaikan terhadap $\mu$.
Dalam kasus pertama, $G\in\Cal G$ dan $G=G\setminus H$ terabaikan terhadap
$\nu$. Dalam kasus kedua, $F_1\setminus G\in\Cal G$ dan
$(F_1\setminus H)\setminus G$ terabaikan terhadap $\nu$. Karena $G$
sembarang, $F_1\setminus H$ adalah atom terhadap $\nu$ yang termuat dalam
$F$; karena $F$ sembarang, $\nu$ atomik murni.

\medskip

{\bf (f)(i)} Misalkan $F\in\Tau$ sedemikian sehingga $B\subseteq F$ dan
$\nu^*B=\nu F$ (132Aa); maka $\phi^{-1}[B]\subseteq\phi^{-1}[F]$, sehingga

\Centerline{$\mu^*\phi^{-1}[B]\le\mu\phi^{-1}[F]=\nu F=\nu^*B$.}

\medskip

\quad{\bf (ii)} $\mu^*A\le\mu^*(\phi^{-1}[\phi[A]])\le\nu^*\phi[A]$
menurut (i).

\medskip

{\bf (g)} Untuk sembarang $W\in\Lambda$,

\Centerline{$\mu(\psi\phi)^{-1}[W]=\mu\phi^{-1}[\psi^{-1}[W]]
=\nu\psi^{-1}[W]=\lambda W$.}
}%end of proof of 234B


\leader{234C}{Ukuran \dvrocolon{citra}}\dvAformerly{1{}12E}\cmmnt{
Konstruksi berikut adalah salah satu cara paling umum bagi ruang-ruang ukur
baru untuk muncul.

\medskip

\noindent}{\bf Proposisi} Misalkan $(X,\Sigma,\mu)$ suatu ruang ukur,
$Y$ sembarang himpunan, dan $\phi:X\to Y$ suatu fungsi. Tetapkan

\Centerline{$\Tau=\{F:F\subseteq Y,\,\phi^{-1}[F]\in\Sigma\}$,
\quad$\nu F=\mu(\phi^{-1}[F])$ untuk setiap $F\in\Tau$.}

\noindent Maka $(Y,\Tau,\nu)$ adalah ruang ukur.

\proof{{\bf (a)} $\emptyset=\phi^{-1}[\emptyset]\in\Sigma$, sehingga
$\emptyset\in\Tau$.

\medskip

{\bf (b)} Jika $F\in\Tau$, maka $\phi^{-1}[F]\in\Sigma$, sehingga
$X\setminus\phi^{-1}[F]\in\Sigma$; tetapi
$X\setminus\phi^{-1}[F]=\phi^{-1}[Y\setminus F]$, sehingga
$Y\setminus F\in\Tau$.

\medskip

{\bf (c)} Jika $\sequencen{F_n}$ suatu barisan di dalam $\Tau$,
maka $\phi^{-1}[F_n]\in\Sigma$ untuk setiap $n$, sehingga
$\bigcup_{n\in\Bbb N}\phi^{-1}[F_n]\in\Sigma$; tetapi
$\phi^{-1}[\bigcup_{n\in\Bbb N}F_n]=\bigcup_{n\in\Bbb N}\phi^{-1}[F_n]$,
sehingga $\bigcup_{n\in\Bbb N}F_n\in\Tau$.

Dengan demikian, $\Tau$ adalah aljabar-sigma.

\medskip

{\bf (d)} $\nu\emptyset=\mu\phi^{-1}[\emptyset]=\mu\emptyset=0$.

\medskip

{\bf (e)} Jika $\sequencen{F_n}$ suatu barisan saling lepas di dalam
$\Tau$, maka $\sequencen{\phi^{-1}[F_n]}$ suatu barisan saling lepas
di dalam $\Sigma$, sehingga

\Centerline{$\nu(\bigcup_{n\in\Bbb N}F_n)
=\mu\phi^{-1}[\bigcup_{n\in\Bbb N}F_n]
=\mu(\bigcup_{n\in\Bbb N}\phi^{-1}[F_n])
=\sum_{n=0}^{\infty}\mu\phi^{-1}[F_n]
=\sum_{n=0}^{\infty}\nu F_n$.}

\noindent Jadi, $\nu$ adalah ukuran.

}%end of proof of 234C

\leader{234D}{Definisi}\dvAformerly{1{}12F}
 Dalam konteks 234C, $\nu$ disebut
{\bf ukuran citra}\cmmnt{ atau {\bf ukuran dorong-maju}};
saya akan menuliskannya sebagai $\mu\phi^{-1}$.

\cmmnt{\medskip

\noindent{\bf Catatan} Barangkali perlu saya katakan bahwa konstruksi ini
tidak selalu menghasilkan secara persis ukuran yang `tepat' pada
$Y$; terdapat keadaan ketika suatu modifikasi dari ukuran
$\mu\phi^{-1}$ yang diuraikan di sini lebih berguna. Namun, saya akan
mencatat keadaan tersebut secara eksplisit saat muncul; ketika saya
menggunakan frasa tanpa hiasan `ukuran citra', yang saya maksud
adalah ukuran yang dibangun di atas.
}%end of comment

\leader{234E}{Proposisi} Misalkan $(X,\Sigma,\mu)$ suatu ruang ukur,
$Y$ suatu himpunan, dan $\phi:X\to Y$ suatu fungsi; misalkan
$\mu\phi^{-1}$ ukuran citra pada $Y$.

(a) $\phi$ bersifat \imp\ untuk $\mu$ dan $\mu\phi^{-1}$.

(b)\dvAformerly{2{}12Bd} Jika $\mu$ lengkap, maka demikian pula
$\mu\phi^{-1}$.

(c)\dvAformerly{1{}12Xd} Jika $Z$ suatu himpunan lain, dan
$\psi:Y\to Z$ suatu fungsi, maka ukuran-ukuran citra
$\mu(\psi\phi)^{-1}$ dan $(\mu\phi^{-1})\psi^{-1}$ pada $Z$ adalah sama.

\proof{{\bf (a)} Langsung dari definisi-definisi.

\medskip

{\bf (b)} Tuliskan $\nu$ untuk $\mu\phi^{-1}$ dan $\Tau$ untuk domainnya.
Jika $\nu^*B=0$, maka $\mu^*\phi^{-1}[B]=0$, menurut 234B(f-i); karena
$\mu$ lengkap, $\phi^{-1}[B]\in\Sigma$, sehingga $B\in\Tau$. Karena $B$
sembarang, $\nu$ lengkap.

\medskip

{\bf (c)} Untuk $G\subseteq Z$ dan $u\in[0,\infty]$,

$$\eqalign{(\mu(\psi\phi)^{-1})(G)&\text{ terdefinisi dan sama dengan }u\cr
&\iff\mu((\psi\phi)^{-1}[G])\text{ terdefinisi dan sama dengan }u\cr
&\iff\mu(\phi^{-1}[\psi^{-1}[G]])\text{ terdefinisi dan sama dengan }u\cr
&\iff(\mu\phi^{-1})(\psi^{-1}[G])\text{ terdefinisi dan sama dengan }u\cr
&\iff((\mu\phi^{-1})\psi^{-1})(G)\text{ terdefinisi dan sama dengan }u.\cr}$$

}%end of proof of 234E


\leader{*234F}{}\dvAformerly{1{}32G}\cmmnt{ Dalam arah sebaliknya,
konstruksi ukuran tarik-balik berikut kadang-kadang berguna.

\medskip

\noindent}{\bf Proposisi} Misalkan $X$ suatu himpunan, $(Y,\Tau,\nu)$
suatu ruang ukur, dan $\phi:X\to Y$ suatu fungsi sedemikian sehingga
$\phi [X]$ mempunyai ukuran luar penuh dalam $Y$. Maka terdapat suatu
ukuran $\mu$ pada $X$, dengan domain
$\Sigma=\{\phi^{-1}[F]:F\in\Tau\}$, sedemikian sehingga $\phi$ bersifat
\imp\ untuk $\mu$ dan $\nu$.

\proof{ Pemeriksaan bahwa $\Sigma$ adalah aljabar-sigma subhimpunan dari
$X$ bersifat langsung; yang perlu kita ketahui hanyalah bahwa
$\phi^{-1}[\emptyset]=\emptyset$,
$X\setminus\phi^{-1}[F]=\phi^{-1}[Y\setminus F]$ untuk setiap
$F\subseteq Y$, dan bahwa
$\phi^{-1}[\bigcup_{n\in\Bbb N}F_n]=\bigcup_{n\in\Bbb N}\phi^{-1}[F_n]$
untuk setiap barisan $\sequencen{F_n}$ dari subhimpunan-subhimpunan $Y$.
Fakta kuncinya ialah bahwa jika $F_1$, $F_2\in\Tau$ dan
$\phi^{-1}[F_1]=\phi^{-1}[F_2]$, maka $\phi[X]$ tidak bertemu dengan
$F_1\symmdiff F_2$; karena $\phi[X]$ mempunyai ukuran luar penuh,
$F_1\symmdiff F_2$ terabaikan terhadap $\nu$ dan $\nu F_1=\nu F_2$.
Dengan demikian, rumus $\mu\phi^{-1}[F]=\nu F$ memang mendefinisikan
suatu fungsi $\mu:\Sigma\to[0,\infty]$. Sekarang

\Centerline{$\mu\emptyset=\mu\phi^{-1}[\emptyset]=\nu\emptyset=0$.}

\noindent Selanjutnya, jika $\sequencen{E_n}$ suatu barisan saling lepas
di dalam $\Sigma$, pilih $F_n\in\Tau$ sedemikian sehingga
$E_n=\phi^{-1}[F_n]$ untuk setiap $n\in\Bbb N$. Barisan
$\sequencen{F_n}$ belum tentu saling lepas, tetapi jika kita menetapkan
$F'_n=F_n\setminus\bigcup_{i<n}F_i$ untuk setiap $n\in\Bbb N$, maka
$\sequencen{F'_n}$ saling lepas dan

\Centerline{$E_n=E_n\setminus\bigcup_{i<n}E_i=\phi^{-1}[F'_n]$}

\noindent untuk setiap $n$; sehingga

\Centerline{$\mu(\bigcup_{n\in\Bbb N}E_n)
=\nu(\bigcup_{n\in\Bbb N}F'_n)
=\sum_{n=0}^{\infty}\nu F'_n
=\sum_{n=0}^{\infty}\mu E_n$.}

\noindent Karena $\sequencen{E_n}$ sembarang, $\mu$ adalah ukuran pada
$X$, sebagaimana diperlukan.
}%end of proof of 234F


\leader{234G}{Jumlah \dvrocolon{ukuran}}\dvAformerly{1{}12Ya;
1{}12Xe untuk jumlah dua ukuran}\cmmnt{ Sekarang saya beralih ke cara
yang sangat berbeda untuk membangun ukuran. Gagasannya jelas, tetapi
perincian teknisnya, dalam kasus umum yang ingin saya kaji, perlu diperhatikan.

\medskip

\noindent}{\bf Proposisi}
Misalkan $X$ suatu himpunan, dan $\familyiI{\mu_i}$ suatu keluarga ukuran
pada $X$. Untuk setiap $i\in I$, misalkan $\Sigma_i$ domain dari $\mu_i$.
Tetapkan $\Sigma=\Cal PX\cap\bigcap_{i\in I}\Sigma_i$ dan definisikan
$\mu:\Sigma\to[0,\infty]$ dengan menetapkan
$\mu E=\sum_{i\in I}\mu_iE$ untuk setiap $E\in\Sigma$. Maka $\mu$
adalah ukuran pada $X$.

\proof{ $\Sigma$ adalah aljabar-sigma subhimpunan dari $X$ karena setiap
$\Sigma_i$ demikian. (Terapkan 111Ga dengan
$\frak S=\{\Sigma_i:i\in I\}\cup\{\Cal PX\}$.)
Tentu saja $\mu$ bernilai di dalam $[0,\infty]$ (226A).
$\mu\emptyset=0$ karena $\mu_i\emptyset=0$ untuk setiap $i$. Jika
$\sequencen{E_n}$ adalah barisan saling lepas di dalam $\Sigma$ dengan
gabungan $E$, maka

$$\eqalignno{\mu E
&=\sum_{i\in I}\mu_iE
=\sum_{i\in I}\sum_{n=0}^{\infty}\mu_iE_n
=\sum_{n=0}^{\infty}\sum_{i\in I}\mu_iE_n\cr
\displaycause{226Af}
&=\sum_{n=0}^{\infty}\mu E_n.\cr}$$

\noindent Jadi, $\mu$ adalah ukuran.
}%end of proof of 234G

\medskip

\noindent{\bf Catatan} Dalam konteks ini, saya akan menyebut $\mu$
{\bf jumlah} dari keluarga $\familyiI{\mu_i}$.

\leader{234H}{Proposisi}\dvAnew{2008}
 Misalkan $X$ suatu himpunan dan $\familyiI{\mu_i}$ suatu keluarga ukuran
lengkap pada $X$ dengan jumlah $\mu$.

(a) $\mu$ lengkap.

(b)(i) Suatu subhimpunan dari $X$ terabaikan terhadap $\mu$ jika dan
hanya jika subhimpunan tersebut terabaikan terhadap $\mu_i$ untuk setiap
$i\in I$.

\quad(ii) Suatu subhimpunan dari $X$ koterabaikan terhadap $\mu$ jika dan
hanya jika subhimpunan tersebut koterabaikan terhadap $\mu_i$ untuk setiap
$i\in I$.

(c)\dvAformerly{2{}12Xh, untuk jumlah dua ukuran}
 Misalkan $f$ suatu fungsi yang didefinisikan pada suatu subhimpunan dari
$X$ dan bernilai dalam $[-\infty,\infty]$. Maka $\int fd\mu$ terdefinisi di
$[-\infty,\infty]$ jika dan hanya jika $\int fd\mu_i$ terdefinisi di
$[-\infty,\infty]$ untuk setiap $i$, dan salah satu dari
$\sum_{i\in I}\int f^+d\mu_i$, $\sum_{i\in I}\int f^-d\mu_i$ berhingga; dalam hal
ini $\int fd\mu=\sum_{i\in I}\int fd\mu_i$.

\proof{ Tuliskan $\Sigma_i=\dom\mu_i$ untuk $i\in I$,
$\Sigma=\Cal PX\cap\bigcap_{i\in I}\Sigma_i=\dom\mu$.

\medskip

{\bf (a)} Jika $E\subseteq F\in\Sigma$ dan $\mu F=0$, maka $\mu_iF=0$
untuk setiap $i\in I$; karena $\mu_i$ lengkap, $E\in\Sigma_i$ untuk
setiap $i\in I$, dan $E\in\Sigma$.

\medskip

{\bf (b)} Sekarang ini langsung mengikuti, karena suatu himpunan
$A\subseteq X$ terabaikan terhadap $\mu$ jika dan hanya jika $\mu A=0$.

\medskip

{\bf (c)(i)} Pertama-tama perhatikan bahwa (b-ii) memberi tahu kita bahwa,
di bawah salah satu hipotesis, $\dom f$ koterabaikan baik terhadap $\mu$
maupun terhadap setiap $\mu_i$. Jadi, jika kita memperluas $f$ ke $X$
dengan memberinya nilai $0$ pada $X\setminus\dom f$, baik $\int fd\mu$
maupun $\sum_{i\in I}\int fd\mu_i$ tidak terpengaruh. Karena itu, mulai
sekarang andaikan bahwa $f$ didefinisikan di seluruh $X$. Sekarang jelas
bahwa salah satu hipotesis memastikan $f$ terukur terhadap $\Sigma$,
yakni terukur terhadap $\Sigma_i$ untuk setiap $i\in I$.

\medskip

\quad{\bf (ii)} Andaikan $f$ nonnegatif.
Untuk $n\in\Bbb N$, tetapkan
$f_n(x)=\sum_{k=1}^{4^n}2^{-n}\chi\{x:f(x)\ge 2^{-n}k\}$, sehingga
$\sequencen{f_n}$ adalah barisan tak-menurun dengan supremum $f$. Kita
mempunyai

$$\eqalign{\int f_nd\mu
&=\sum_{k=1}^{4^n}2^{-n}\mu\{x:f(x)\ge 2^{-n}k\}
=\sum_{k=1}^{4^n}\sum_{i\in I}2^{-n}\mu_i\{x:f(x)\ge 2^{-n}k\}\cr
&=\sum_{i\in I}\sum_{k=1}^{4^n}2^{-n}\mu_i\{x:f(x)\ge 2^{-n}k\}
=\sum_{i\in I}\int f_nd\mu_i\cr}$$

\noindent untuk setiap $n$, sehingga

$$\eqalign{\int fd\mu
&=\sup_{n\in\Bbb N}\int f_nd\mu
=\sup_{n\in\Bbb N}\sup_{J\subseteq I\text{ berhingga}}
   \sum_{i\in J}\int f_nd\mu_i
=\sup_{J\subseteq I\text{ berhingga}}\sup_{n\in\Bbb N}
   \sum_{i\in J}\int f_nd\mu_i\cr
&=\sup_{J\subseteq I\text{ berhingga}}
   \lim_{n\to\infty}\sum_{i\in J}\int f_nd\mu_i
=\sup_{J\subseteq I\text{ berhingga}}
   \sum_{i\in J}\lim_{n\to\infty}\int f_nd\mu_i
=\sum_{i\in I}\int fd\mu_i.\cr}$$

\medskip

\quad{\bf (iii)} Secara umum,

$$\eqalign{\int fd\mu&\text{ terdefinisi di }[-\infty,\infty]\cr
&\iff\int f^+d\mu\text{ dan }\int f^-d\mu
  \text{ terdefinisi dan paling banyak salah satunya tak berhingga}\cr
&\iff\sum_{i\in I}\int f^+d\mu_i\text{ dan }\sum_{i\in I}\int f^-d\mu_i
  \text{ terdefinisi dan paling banyak salah satunya tak berhingga}\cr
&\iff\int fd\mu_i\text{ terdefinisi untuk setiap }i
  \text{ dan paling banyak salah satu dari }\sum_{i\in I}\int f^+d\mu_i,\cr
&\mskip400mu
  \sum_{i\in I}\int f^-d\mu_i\text{ tak berhingga},\cr}$$

\noindent dan dalam hal ini

\Centerline{$\int fd\mu
=\int f^+d\mu-\int f^-d\mu
=\sum_{i\in I}\int f^+d\mu_i-\sum_{i\in I}\int f^-d\mu_i
=\sum_{i\in I}\int fd\mu_i$.}
}%end of proof of 234H

\leader{234I}{Ukuran
\dvrocolon{integral tak tentu}}\dvAformerly{2{}34A}\cmmnt{ Dengan memperluas
gagasan yang telah digunakan dalam 232D, kita sampai pada konstruksi berikut;
sekali lagi, kita perlu berhati-hati terhadap perincian formal jika ingin
memperoleh manfaat sepenuhnya.

\medskip

\noindent}{\bf Teorema} Misalkan $(X,\Sigma,\mu)$ suatu ruang ukur,
dan $f$ suatu fungsi bernilai real nonnegatif yang terukur secara virtual
terhadap $\mu$ dan didefinisikan pada suatu subhimpunan koterabaikan dari
$X$. Tuliskan $\nu F=\int f\times\chi F\,d\mu$ kapan pun $F\subseteq X$
sedemikian sehingga integral tersebut terdefinisi di
$[0,\infty]$\cmmnt{ menurut konvensi 133A}.
Maka $\nu$ adalah ukuran lengkap pada $X$, dan domainnya mencakup
$\Sigma$.
%the trouble with taking [0,\infty]-valued functions is that if
%I allow  f=\infty  on a non-negligible set then  \nu  becomes
%non-semi-finite for a silly reason

\proof{{\bf (a)} Tuliskan $\Tau$ untuk domain $\nu$, yakni keluarga
himpunan $F\subseteq X$ sedemikian sehingga
$\int f\times\chi F\,d\mu$ terdefinisi di $[0,\infty]$, yakni
$f\times\chi F$ terukur secara virtual terhadap $\mu$ (133A). Maka
$\Tau$ adalah aljabar-sigma subhimpunan dari $X$. \Prf\
Untuk setiap $F\in\Tau$, misalkan $H_F\subseteq X$ suatu himpunan
koterabaikan terhadap $\mu$ sedemikian sehingga
$f\times\chi F\restr H_F$ terukur terhadap $\Sigma$. Karena $f$ sendiri
terukur secara virtual terhadap $\mu$, $X\in\Tau$. Jika $F\in\Tau$, maka

\Centerline{$f\times\chi(X\setminus F)\restr(H_X\cap H_F)
=f\restr(H_X\cap H_F)-(f\times\chi F)\restr(H_X\cap H_F)$}

\noindent terukur terhadap $\Sigma$, sedangkan $H_X\cap H_F$
koterabaikan terhadap $\mu$, sehingga $X\setminus F\in\Tau$. Jika
$\sequencen{F_n}$ adalah barisan di dalam $\Tau$ dengan gabungan $F$,
tetapkan $H=\bigcap_{n\in\Bbb N}H_{F_n}$; maka $H$ koterabaikan,
$f\times\chi F_n\restr H$ terukur terhadap $\Sigma$ untuk setiap
$n\in\Bbb N$, dan $f\times\chi F=\sup_{n\in\Bbb N}f\times\chi F_n$,
sehingga $f\times\chi F\restr H$ terukur terhadap $\Sigma$, dan
$F\in\Tau$. Dengan demikian, $\Tau$ adalah aljabar-sigma. Jika
$F\in\Sigma$, maka $f\times\chi F\restr H_X$ terukur terhadap $\Sigma$,
sehingga $F\in\Tau$.\ \Qed

\medskip

{\bf (b)} Selanjutnya, $\nu$ adalah ukuran. \Prf\ Tentu saja
$\nu F\in[0,\infty]$ untuk setiap $F\in\Tau$.
$f\times\chi\emptyset=0$ di mana pun fungsi tersebut didefinisikan,
sehingga $\nu\emptyset=0$. Jika $\sequencen{F_n}$ adalah barisan saling
lepas di dalam $\Tau$ dengan gabungan $F$, maka
$f\times\chi F=\sum_{n=0}^{\infty}f\times\chi F_n$.
Jika $\nu F_m=\infty$ untuk suatu $m$, maka tentu saja
$\nu F=\infty=\sum_{n=0}^{\infty}\nu F_n$. Jika $\nu F_m<\infty$ untuk
setiap $m$, tetapi $\sum_{n=0}^{\infty}\nu F_n=\infty$, maka

\Centerline{$\int f\times\chi(\bigcup_{n\le m}F_n)
=\sum_{n=0}^m\int f\times\chi F_n\to\infty$}

\noindent ketika $m\to\infty$, sehingga sekali lagi
$\nu F=\infty=\sum_{n=0}^{\infty}\nu F_n$. Jika
$\sum_{n=0}^{\infty}\nu F_n<\infty$, maka menurut teorema B.Levi

\Centerline{$\nu F=\int\sum_{n=0}^{\infty}f\times\chi F_n
=\sum_{n=0}^{\infty}\int f\times\chi F_n
=\sum_{n=0}^{\infty}\nu F_n$.
\Qed}

\medskip

{\bf (c)} Terakhir, $\nu$ lengkap. \Prf\ Jika $A\subseteq F\in\Tau$
dan $\nu F=0$, maka $f\times\chi F=0$ hampir di mana-mana, sehingga
$f\times\chi A=0$ hampir di mana-mana dan $\nu A$ terdefinisi serta sama
dengan nol.\ \Qed
}%end of proof of 234I

\leader{234J}{Definisi}\dvAformerly{2{}34B}
Misalkan $(X,\Sigma,\mu)$ suatu ruang ukur, dan
$\nu$ ukuran lain pada $X$ dengan domain $\Tau$. Saya akan menyebut
$\nu$ suatu {\bf ukuran integral tak tentu} terhadap $\mu$, atau
kadang-kadang suatu {\bf ukuran integral tak tentu yang dilengkapi}, jika
ukuran tersebut dapat diperoleh dengan metode 234I dari suatu fungsi
nonnegatif $f$ yang terukur secara virtual dan didefinisikan hampir di
mana-mana dalam $X$. Dalam hal ini, $f$ adalah suatu turunan
Radon-Nikod\'ym dari $\nu$ terhadap $\mu$ dalam pengertian 232Hf.
\cmmnt{Seperti dalam 232Hf, frasa {\bf fungsi kerapatan} juga digunakan
dalam konteks ini.}

\leader{234K}{Catatan-catatan}\dvAformerly{2{}34C}
Misalkan $(X,\Sigma,\mu)$ suatu ruang ukur, dan $f$ suatu fungsi bernilai
real nonnegatif yang terukur secara virtual terhadap $\mu$ dan
didefinisikan hampir di mana-mana dalam $X$; misalkan $\nu$ ukuran
integral tak tentu yang bersesuaian.

\spheader 234Ka\cmmnt{ Terdapat suatu fungsi yang terukur terhadap
$\Sigma$, $g:X\to\coint{0,\infty}$, sedemikian sehingga
$f=g\,\,\mu$-hampir di mana-mana. \Prf\ Misalkan
$H\subseteq\dom f$ suatu himpunan terukur yang koterabaikan sedemikian
sehingga $f\restr H$ terukur, dan tetapkan $g(x)=f(x)$ untuk $x\in H$,
$g(x)=0$ untuk $x\in X\setminus H$.\ \QeD\ Dalam hal ini,
$\int f\times\chi E\,d\mu=\int g\times\chi E\,d\mu$ jika salah satunya
terdefinisi. Jadi, $g$ adalah suatu turunan Radon-Nikod\'ym dari $\nu$, dan}
$\nu$ mempunyai suatu turunan Radon-Nikod\'ym yang terukur terhadap
$\Sigma$ dan didefinisikan di mana-mana.

%includes old 2{}34Ce
\spheader 234Kb Jika $E$ terabaikan terhadap $\mu$,
maka\cmmnt{ $f\times\chi E=0\,\,\mu$-hampir di mana-mana, sehingga}
$\nu E=0$.
\cmmnt{Banyak penulis bersedia mengatakan `$\nu$ kontinu mutlak
terhadap $\mu$' dalam konteks ini. Namun, jika $\nu$ tidak berhingga total,
ia belum tentu kontinu mutlak dalam pengertian
$\epsilon$-$\delta$ pada 232Aa (234Xh), dan kesulitan lebih lanjut dapat
muncul jika $\mu$ atau $\nu$ tidak $\sigma$-hingga (lihat 234Yk, 234Ym).
}%end of comment

\cmmnt{\spheader 234Kc Saya telah mendefinisikan `ukuran integral
tak tentu' sedemikian rupa sehingga menghasilkan ukuran lengkap. Menurut
saya, inilah yang paling masuk akal dalam kebanyakan penerapan. Ada
kesempatan ketika tampaknya lebih tepat menggunakan ukuran
$\nu_0:\Sigma\to[0,\infty]$ yang didefinisikan dengan menetapkan
$\nu_0 E=\int_Efd\mu=\int f\times\chi E\,d\mu$ untuk $E\in\Sigma$.
Saya kira saya akan menyebut ini {\bf ukuran integral tak tentu yang belum
dilengkapi} terhadap $\mu$ yang didefinisikan oleh $f$. ($\nu$ selalu
merupakan pelengkapan dari $\nu_0$; lihat 234Lb.)
}%end of comment

\spheader 234Kd\cmmnt{ Perhatikan cara saya merumuskan definisi $\nu$:
`$\nu E=\int f\times\chi E\,d\mu$ jika integralnya terdefinisi',
alih-alih `$\nu E=\int_Efd\mu$'. Pokoknya ialah bahwa rumus yang
lebih panjang memberikan kaidah untuk menentukan domain $\nu$. Tentu saja
memang benar bahwa} $\nu E=\int_Efd\mu$ untuk setiap
$E\in\dom\nu$\cmmnt{ (terapkan 214F pada $f\times\chi E$)}.

\spheader 234Ke\dvAnew{2008}
\cmmnt{Karena }$\mu$ dan pelengkapannya mendefinisikan
\cmmnt{fungsi-fungsi terukur secara virtual yang sama, ideal nol yang sama,
dan integral-integral yang sama\cmmnt{ (212Eb, 212F)}, keduanya mendefinisikan} ukuran-ukuran
integral tak tentu yang sama.

\leader{234L}{Domain suatu ukuran
\dvrocolon{integral tak tentu}}\dvAformerly{2{}34D}\cmmnt{ Kadang-kadang
berguna untuk memiliki uraian eksplisit mengenai domain suatu ukuran yang
dibangun dengan cara ini.

\medskip

\noindent}{\bf Proposisi} Misalkan $(X,\Sigma,\mu)$ suatu ruang ukur,
$f$ suatu fungsi nonnegatif yang terukur secara virtual terhadap $\mu$
dan didefinisikan hampir di mana-mana dalam $X$, dan $\nu$ ukuran
integral tak tentu yang bersesuaian. Tetapkan
$G=\{x:x\in\dom f$, $f(x)>0\}$, dan misalkan $\hat\Sigma$ domain
pelengkapan $\hat\mu$ dari $\mu$.

(a) Domain $\Tau$ dari $\nu$ adalah
$\{E:E\subseteq X,\,E\cap G\in\hat\Sigma\}$; \cmmnt{khususnya,}
$\Tau\supseteq\hat\Sigma\supseteq\Sigma$.

(b) $\nu$ adalah pelengkapan dari pembatasannya pada $\Sigma$.

(c) Suatu himpunan $A\subseteq X$ terabaikan terhadap $\nu$ jika dan
hanya jika $A\cap G$ terabaikan terhadap $\mu$.

(d) Khususnya, jika $\mu$\cmmnt{ sendiri} lengkap, maka
$\Tau=\{E:E\subseteq X,\,E\cap G\in\Sigma\}$ dan $\nu A=0$ jika dan
hanya jika $\mu(A\cap G)=0$.

\proof{{\bf (a)(i)} Jika $E\in\Tau$, maka $f\times\chi E$ terukur secara
virtual, sehingga terdapat suatu himpunan terukur koterabaikan
$H\subseteq\dom f$ sedemikian sehingga $f\times\chi E\restr H$ terukur.
Sekarang $E\cap G\cap H=\{x:x\in H,\,(f\times\chi E)(x)>0\}$ harus
termasuk dalam $\Sigma$, sedangkan $E\cap G\setminus H$ terabaikan,
sehingga termasuk dalam $\hat\Sigma$, dan $E\cap G\in\hat\Sigma$.

\medskip

\quad{\bf (ii)} Jika $E\cap G\in\hat\Sigma$, misalkan
$F_1$, $F_2\in\Sigma$ sedemikian sehingga
$F_1\subseteq E\cap G\subseteq F_2$ dan $F_2\setminus F_1$ terabaikan.
Misalkan $H\subseteq\dom f$ suatu himpunan koterabaikan sedemikian sehingga
$f\restr H$ terukur. Maka $H'=H\setminus(F_2\setminus F_1)$ koterabaikan
dan $f\times\chi E\restr H'=f\times\chi F_1\restr H'$ terukur, sehingga
$f\times\chi E$ terukur secara virtual dan $E\in\Tau$.

\medskip

{\bf (b)} Jadi, rumus yang diberikan memang menguraikan $\Tau$. Jika
$E\in\Tau$, misalkan $F_1$, $F_2\in\Sigma$ sedemikian sehingga
$F_1\subseteq E\cap G\subseteq F_2$ dan $\mu(F_2\setminus F_1)=0$.
Karena $G$ sendiri juga termasuk dalam $\hat\Sigma$, terdapat
$G_1$, $G_2\in\Sigma$ sedemikian sehingga
$G_1\subseteq G\subseteq G_2$ dan $\mu(G_2\setminus G_1)=0$. Tetapkan
$F_2'=F_2\cup(X\setminus G_1)$. Maka $F_2'\in\Sigma$ dan
$F_1\subseteq E\subseteq F'_2$. Tetapi
$(F'_2\setminus F_1)\cap G
\subseteq(G_2\setminus G_1)\cup(F_2\setminus F_1)$ terabaikan terhadap
$\mu$, sehingga $\nu(F'_2\setminus F_1)=0$.

Ini menunjukkan bahwa jika $\nuprime$ adalah pelengkapan dari
$\nu\restr\Sigma$ dan $\Tau'$ domainnya, maka $\Tau\subseteq\Tau'$.
Namun, karena $\nu$ lengkap, tentu saja ia memperluas $\nuprime$, sehingga
$\nu=\nuprime$, sebagaimana dinyatakan.

\medskip

{\bf (c)} Sekarang ambil sembarang $A\subseteq X$. Karena $\nu$ lengkap,

$$\eqalign{A\text{ terabaikan terhadap }\nu\text{}
&\iff\nu A=0\cr
&\iff\int f\times\chi A\,d\mu=0\cr
&\iff f\times\chi A=0\,\,\mu\text{-hampir di mana-mana}\cr
&\iff A\cap G\text{ terabaikan terhadap }\mu\text{.}\cr}$$

\medskip

{\bf (d)} Ini hanyalah pernyataan ulang dari (a) dan (c) ketika
$\mu=\hat\mu$.
}%end of proof of 234L

\leader{234M}{Akibat}\dvAformerly{2{}34E}
 Jika $(X,\Sigma,\mu)$ suatu ruang ukur lengkap dan $G\in\Sigma$, maka
ukuran integral tak tentu terhadap $\mu$ yang didefinisikan oleh $\chi G$
hanyalah ukuran $\mu\LLcorner G$ yang didefinisikan dengan menetapkan

\Centerline{$(\mu\LLcorner G)(F)=\mu(F\cap G)$ kapan pun $F\subseteq X$
dan $F\cap G\in\Sigma$.}

\proof{ 234Ld.
}%end of proof of 234M

\leader{*234N}{}\cmmnt{ Dua hasil berikut tidak akan diandalkan dalam
jilid ini, tetapi saya menyertakannya sebagai acuan di masa depan dan
untuk memberikan gambaran mengenai jangkauan gagasan-gagasan ini.

\medskip

\noindent}{\bf Proposisi}\dvAformerly{2{}34F}
Misalkan $(X,\Sigma,\mu)$ suatu ruang ukur, dan
$\nu$ suatu ukuran integral tak tentu terhadap $\mu$.

(a) Jika $\mu$ semihingga, maka demikian pula $\nu$.

(b) Jika $\mu$ lengkap dan ditentukan secara lokal, maka demikian pula
$\nu$.

(c) Jika $\mu$ dapat dilokalkan, maka demikian pula $\nu$.

(d) Jika $\mu$ dapat dilokalkan secara ketat, maka demikian pula $\nu$.

(e) Jika $\mu$ $\sigma$-hingga, maka demikian pula $\nu$.

(f) Jika $\mu$ tanpa atom, maka demikian pula $\nu$.

\proof{ Menurut 234Ka, kita dapat menyatakan $\nu$ sebagai integral tak
tentu dari suatu fungsi yang terukur terhadap $\Sigma$,
$f:X\to\coint{0,\infty}$. Misalkan $\Tau$ domain $\nu$, dan
$\hat\Sigma$ domain pelengkapan $\hat\mu$ dari $\mu$; tetapkan
$G=\{x:x\in X,\,f(x)>0\}\in\Sigma$.

\medskip

{\bf (a)} Andaikan $E\in\Tau$ dan $\nu E=\infty$. Maka $E\cap G$ tidak
mungkin terabaikan terhadap $\mu$. Karena $\mu$ semihingga, terdapat
suatu $F\in\Sigma$ yang tidak terabaikan sedemikian sehingga
$F\subseteq E\cap G$ dan $\mu F<\infty$. Sekarang
$F=\bigcup_{n\in\Bbb N}\{x:x\in F$, $2^{-n}\le f(x)\le n\}$, sehingga
terdapat suatu $n\in\Bbb N$ sedemikian sehingga
$F'=\{x:x\in F$, $2^{-n}\le f(x)\le n\}$ tidak terabaikan.
Karena $f$ terukur, $F'\in\Sigma\subseteq\Tau$ dan
$2^{-n}\mu F'\le\nu F'\le n\mu F'$. Dengan demikian, kita telah menemukan
suatu $F'\subseteq E$ sedemikian sehingga $0<\nu F'<\infty$. Karena $E$
sembarang, $\nu$ semihingga.

\medskip

{\bf (b)} Kita telah mengetahui bahwa $\nu$ lengkap (234Lb) dan
semihingga. Sekarang andaikan $E\subseteq X$ sedemikian sehingga
$E\cap F\in\Tau$, yakni $E\cap F\cap G\in\Sigma$ (234Ld), kapan pun
$F\in\Tau$ dan $\nu F<\infty$. Maka $E\cap G\cap F\in\Sigma$ kapan pun
$F\in\Sigma$ dan $\mu F<\infty$. \Prf\ Tetapkan
$F_n=\{x:x\in F\cap G,\,f(x)\le n\}$. Maka
$\nu F_n\le n\mu F<\infty$, sehingga $E\cap G\cap F_n\in\Sigma$ untuk
setiap $n$. Namun, ini berarti
$E\cap G\cap F=\bigcup_{n\in\Bbb N}E\cap G\cap F_n\in\Sigma$.\ \QeD\ 
Karena $\mu$ ditentukan secara lokal, $E\cap G\in\Sigma$ dan $E\in\Tau$.
Karena $E$ sembarang, $\nu$ ditentukan secara lokal.

\medskip

{\bf (c)} Misalkan $\Cal F\subseteq\Tau$ sembarang himpunan. Tetapkan
$\Cal E=\{F\cap G:F\in\Cal F\}$, sehingga
$\Cal E\subseteq\hat\Sigma$. Menurut 212Ga, $\hat\mu$ dapat dilokalkan,
sehingga $\Cal E$ mempunyai suatu supremum esensial $H\in\hat\Sigma$.
Namun sekarang, untuk sembarang $H'\in\Tau$,
$H'\cup(X\setminus G)=(H'\cap G)\cup(X\setminus G)$ termasuk dalam
$\hat\Sigma$, sehingga

$$\eqalign{\nu(F\setminus H')&=0\text{ untuk setiap }F\in\Cal F\cr
&\iff\hat\mu(F\cap G\setminus H')=0\text{ untuk setiap }F\in\Cal F\cr
&\iff\hat\mu(E\setminus H')=0\text{ untuk setiap }E\in\Cal E\cr
&\iff\hat\mu(E\setminus(H'\cup(X\setminus G)))=0
  \text{ untuk setiap }E\in\Cal E\cr
&\iff\hat\mu(H\setminus(H'\cup(X\setminus G)))=0\cr
&\iff\hat\mu(H\cap G\setminus H')=0\cr
&\iff\nu(H\setminus H')=0.\cr}$$

\noindent Jadi, $H$ juga merupakan supremum esensial dari $\Cal F$
di dalam $\Tau$. Karena $\Cal F$ sembarang, $\nu$ dapat dilokalkan.

\medskip

{\bf (d)} Misalkan $\langle X_i\rangle_{i\in I}$ suatu dekomposisi
$X$ untuk $\mu$; maka keluarga tersebut juga merupakan dekomposisi untuk
$\hat\mu$ (212Gb). Tetapkan $F_0=X\setminus G$,
$F_n=\{x:x\in G,\,n-1<f(x)\le n\}$ untuk $n\ge 1$. Maka
$\langle X_i\cap F_n\rangle_{i\in I,n\in\Bbb N}$ adalah dekomposisi
untuk $\nu$. \Prf\ (i) $\langle X_i\rangle_{i\in I}$ dan
$\langle F_n\rangle_{n\in\Bbb N}$ adalah partisi $X$ menjadi
anggota-anggota $\Sigma\subseteq\Tau$, sehingga
$\langle X_i\cap F_n\rangle_{i\in I,n\in\Bbb N}$ juga demikian.
(ii) $\nu(X_i\cap F_0)=0$,
$\nu(X_i\cap F_n)\le n\mu X_i<\infty$ untuk $i\in I$, $n\ge 1$.
(iii) Jika $E\subseteq X$ dan $E\cap X_i\cap F_n\in\Tau$ untuk setiap
$i\in I$ dan $n\in\Bbb N$, maka
$E\cap X_i\cap G=\bigcup_{n\in\Bbb N}E\cap X_i\cap F_n\cap G$ termasuk
dalam $\hat\Sigma$ untuk setiap $i$, sehingga
$E\cap G\in\hat\Sigma$ dan $E\in\Tau$.
(iv) Jika $E\in\Tau$, maka tentu saja

\Centerline{$\sum_{i\in I,n\in\Bbb N}\nu(E\cap X_i\cap F_n)
=\sup_{J\subseteq I\times\Bbb N\text{ berhingga}}
  \sum_{(i,n)\in J}\nu(E\cap X_i\cap F_n)\le\nu E$.}

\noindent Jadi, jika
$\sum_{i\in I,n\in\Bbb N}\nu(E\cap X_i\cap F_n)=\infty$, jumlah tersebut
tentu sama dengan $\nu E$. Jika jumlahnya berhingga, maka
$K=\{i:i\in I,\,\nu(E\cap X_i)>0\}$ harus terhitung. Namun, untuk
$i\in I\setminus K$, $\int_{E\cap X_i}fd\mu=0$, sehingga
$f=0\,\,\mu$-hampir di mana-mana\ pada $E\cap X_i$, yakni
$\hat\mu(E\cap G\cap X_i)=0$. Karena
$\langle X_i\rangle_{i\in I}$ merupakan dekomposisi untuk $\hat\mu$,
$\hat\mu(E\cap G\cap\bigcup_{i\in I\setminus K}X_i)=0$ dan
$\nu(E\cap\bigcup_{i\in I\setminus K}X_i)=0$. Namun, ini berarti bahwa

\Centerline{$\nu E=\sum_{i\in K}\nu(E\cap X_i)
=\sum_{i\in K,n\in\Bbb N}\nu(E\cap X_i\cap F_n)
=\sum_{i\in I,n\in\Bbb N}\nu(E\cap X_i\cap F_n)$.}

\noindent Karena $E$ sembarang,
$\langle X_i\cap F_n\rangle_{i\in I,n\in\Bbb N}$ adalah dekomposisi
untuk $\nu$.\ \QeD\ Jadi, $\nu$ dapat dilokalkan secara ketat.

\medskip

{\bf (e)} Jika $\mu$ $\sigma$-hingga, maka dalam (d) kita dapat mengambil
$I$ terhitung, sehingga $I\times\Bbb N$ juga terhitung, dan $\nu$ akan
$\sigma$-hingga.

\medskip

{\bf (f)} Jika $\mu$ tanpa atom, demikian pula $\hat\mu$ (212Gd).
Jika $E\in\Tau$ dan $\nu E>0$, maka $\hat\mu(E\cap G)>0$, sehingga
terdapat suatu $F\in\hat\Sigma$ sedemikian sehingga
$F\subseteq E\cap G$ dan baik $F$ maupun $E\cap G\setminus F$ tidak
terabaikan terhadap $\hat\mu$. Namun, dalam hal ini, baik
$\nu F=\int_Ffd\mu$ maupun
$\nu(E\setminus F)=\int_{E\setminus F}fd\mu$ pasti lebih besar daripada
$0$ (122Rc). Karena $E$ sembarang, $\nu$ tanpa atom.
}%end of proof of 234N

\leader{*234O}{}\cmmnt{ Untuk ukuran yang dapat dilokalkan, terdapat
uraian langsung mengenai ukuran integral tak tentu yang bersesuaian.

\medskip

\noindent}{\bf Teorema}\dvAformerly{2{}34G}
 Misalkan $(X,\Sigma,\mu)$ suatu ruang ukur yang dapat dilokalkan. Maka
suatu ukuran $\nu$, dengan domain $\Tau\supseteq\Sigma$, adalah ukuran
integral tak tentu terhadap $\mu$ jika dan hanya jika
($\alpha$) $\nu$ semihingga dan nol pada himpunan-himpunan yang terabaikan
terhadap $\mu$
($\beta$) $\nu$ adalah pelengkapan dari pembatasannya pada $\Sigma$
($\gamma$) kapan pun $\nu E>0$, terdapat suatu $F\subseteq E$ sedemikian
sehingga $F\in\Sigma$, $\mu F<\infty$, dan $\nu F>0$.

\proof{{\bf (a)} Jika $\nu$ suatu ukuran integral tak tentu terhadap
$\mu$, maka menurut 234Na, 234Kb, dan 234Lb, ukuran tersebut semihingga,
nol pada himpunan-himpunan yang terabaikan terhadap $\mu$, dan merupakan
pelengkapan dari pembatasannya pada $\Sigma$. Sekarang andaikan
$E\in\Tau$ dan $\nu E>0$. Maka terdapat suatu
$E_0\in\Sigma$ sedemikian sehingga $E_0\subseteq E$ dan
$\nu E_0=\nu E$, menurut 234Lb. Jika $f:X\to\Bbb R$ adalah suatu fungsi
yang terukur terhadap $\Sigma$ dan merupakan turunan Radon-Nikod\'ym dari
$\nu$ (234Ka), dan
$G=\{x:f(x)>0\}$, maka $\mu(G\cap E_0)>0$; karena $\mu$ semihingga,
terdapat suatu $F\in\Sigma$ sedemikian sehingga
$F\subseteq G\cap E_0$ dan $0<\mu F<\infty$, dan dalam hal ini
$\nu F>0$.

\medskip

{\bf (b)} Jadi, sekarang andaikan $\nu$ memenuhi syarat-syarat tersebut.

\medskip

\quad{\bf (i)} Tetapkan
$\Cal E=\{E:E\in\Sigma,\,\nu E<\infty\}$. Untuk setiap $E\in\Cal E$,
tinjau $\nu_E:\Sigma\to\Bbb R$, dengan menetapkan
$\nu_EG=\nu(G\cap E)$ untuk setiap $G\in\Sigma$. Maka $\nu_E$ aditif
terhitung dan kontinu sejati terhadap $\mu$. \Prf\
$\nu_E$ aditif terhitung, persis seperti dalam 231De. Karena $\nu$ nol
pada himpunan-himpunan yang terabaikan terhadap $\mu$, $\nu_E$ harus
kontinu mutlak terhadap $\mu$, menurut 232Ba. Karena $\nu_E$ jelas
memenuhi syarat ($\gamma$) dari 232Bb, ia harus kontinu sejati.\ \Qed

Menurut 232E, terdapat suatu fungsi yang terintegralkan terhadap $\mu$,
$f_E$, sedemikian sehingga $\nu_EG=\int_Gf_Ed\mu$ untuk setiap
$G\in\Sigma$, dan kita dapat mengandaikan bahwa $f_E$ terukur terhadap
$\Sigma$ (232He). Karena $\nu_E$ nonnegatif,
$f_E\ge 0\,\,\mu$-hampir di mana-mana.

\medskip

\quad{\bf (ii)} Jika $E$, $F\in\Cal E$, maka
$f_E=f_F\,\,\mu$-hampir di mana-mana pada $E\cap F$, karena

\Centerline{$\int_Gf_Ed\mu=\nu G=\int_Gf_Fd\mu$}

\noindent kapan pun $G\in\Sigma$ dan $G\subseteq E\cap F$. Karena
$(X,\Sigma,\mu)$ dapat dilokalkan, terdapat suatu fungsi terukur
$f:X\to\Bbb R$ sedemikian sehingga $f_E=f\,\,\mu$-hampir di mana-mana\ 
pada $E$ untuk setiap $E\in\Cal E$ (213N). Karena setiap $f_E$
nonnegatif hampir di mana-mana, kita dapat mengandaikan bahwa $f$
nonnegatif, sebab tentu saja $f_E=f\vee\tbf{0}\,\,\mu$-hampir di
mana-mana pada $E$ untuk setiap $E\in\Cal E$.

\medskip

\quad{\bf (iii)} Misalkan $\nuprime$ ukuran integral tak tentu yang
didefinisikan oleh $f$. Jika $E\in\Cal E$, maka

\Centerline{$\nu E=\int_Ef_Ed\mu=\int_Efd\mu=\nuprime E$.}
\noindent Untuk $E\in\Sigma\setminus\Cal E$, kita mempunyai

\Centerline{$\nuprime E
\ge\sup\{\nuprime F:F\in\Cal E,\,F\subseteq E\}
=\sup\{\nu F:F\in\Cal E,\,F\subseteq E\}=\nu E=\infty$}

\noindent karena $\nu$ semihingga. Jadi, $\nuprime$ dan $\nu$ berimpit
pada $\Sigma$. Namun, karena masing-masing merupakan pelengkapan dari
pembatasannya pada $\Sigma$, keduanya harus sama.
}%end of proof of 234O

\leader{234P}{Pengurutan \dvrocolon{ukuran}}\cmmnt{ Terdapat banyak cara
suatu ukuran dapat mendominasi ukuran lain. Di sini saya akan menguraikan
salah satu cara yang paling sederhana.

\medskip

\noindent}{\bf Definisi}\dvAnew{2008}
Misalkan $\mu$, $\nu$ dua ukuran pada suatu himpunan $X$.
Saya akan mengatakan bahwa $\mu\le\nu$ jika $\mu E$ terdefinisi dan
$\mu E\le\nu E$ kapan pun $\nu$ mengukur $E$.

\leader{234Q}{Proposisi}\dvAnew{2008}
Misalkan $X$ suatu himpunan, dan tuliskan $\Mu$ untuk himpunan semua
ukuran pada $X$.

(a) Dengan mendefinisikan $\le$ seperti dalam 234P, $(\Mu,\le)$ adalah
suatu himpunan terurut parsial.

(b) Jika $\mu$, $\nu\in\Mu$, maka $\mu\le\nu$ jika dan hanya jika terdapat
suatu $\lambda\in\Mu$ sedemikian sehingga $\mu+\lambda=\nu$.

(c) Jika $\mu\le\nu$ dalam $\Mu$ dan $f$ adalah suatu fungsi bernilai
$[-\infty,\infty]$ yang didefinisikan pada suatu subhimpunan dari $X$,
sedemikian sehingga $\int fd\nu$ terdefinisi di $[-\infty,\infty]$, maka
$\int fd\mu$ terdefinisi; jika $f$ nonnegatif,
$\int fd\mu\le\int fd\nu$.

\proof{{\bf (a)} Tentu saja $\mu\le\mu$ untuk setiap $\mu\in\Mu$.
Jika $\mu\le\nu$ dan $\nu\le\lambda$ dalam $\Mu$, maka
$\dom\lambda\subseteq\dom\nu\subseteq\dom\mu$, dan
$\mu E\le\nu E\le\lambda E$ kapan pun $\lambda$ mengukur $E$.
Jika $\mu\le\nu$ dan $\nu\le\mu$, maka
$\dom\mu\subseteq\dom\nu\subseteq\dom\mu$ dan
$\mu E\le\nu E\le\mu E$ untuk setiap $E$ dalam domain bersama keduanya,
sehingga $\mu=\nu$.

\medskip

{\bf (b)(i)}
Jika $\mu+\lambda=\nu$, maka definisi-definisi dalam 234G dan 234P
langsung menunjukkan bahwa $\mu\le\nu$.

\medskip

\quad{\bf (ii)}\grheada\ Dalam arah sebaliknya, jika $\mu\le\nu$,
tuliskan $\Tau$ untuk domain $\nu$.
Definisikan $\lambda:\Tau\to[0,\infty]$ dengan menetapkan

\Centerline{$\lambda G
=\sup\{\nu F-\mu F:F\in\Tau,\,F\subseteq G,\,\mu F<\infty\}$}

\noindent untuk $G\in\Tau$.
Maka $\lambda\in\Mu$. \Prf\ Tentu saja $\dom\lambda=\Tau$ adalah
aljabar-sigma, dan $\lambda\emptyset=0$. Andaikan
$\sequencen{G_n}$ adalah suatu barisan saling lepas di dalam $\Tau$
dengan gabungan $G$. Jika $F\in\Tau$, $F\subseteq G$, dan
$\mu F<\infty$, maka

$$\eqalign{\nu F-\mu F
&=\sum_{n=0}^{\infty}\nu(F\cap G_n)-\sum_{n=0}^{\infty}\mu(F\cap G_n)\cr
&=\sum_{n=0}^{\infty}(\nu(F\cap G_n)-\mu(F\cap G_n))
\le\sum_{n=0}^{\infty}\lambda G_n;\cr}$$

\noindent karena $F$ sembarang,
$\lambda G\le\sum_{n=0}^{\infty}\lambda G_n$. Jika
$\gamma<\sum_{n=0}^{\infty}\lambda G_n$, terdapat suatu $m\in\Bbb N$
sedemikian sehingga $\gamma<\sum_{n=0}^m\lambda G_n$, dan terdapat
$F_0,\ldots,F_m$ sedemikian sehingga $F_n\in\Tau$, $F_n\subseteq G_n$,
dan $\mu F_n<\infty$ untuk setiap $n\le m$, sedangkan
$\sum_{n=0}^m(\nu F_n-\mu F_n)\ge\gamma$. Tetapkan
$F=\bigcup_{n\le m}F_n$; maka $F\in\Tau$, $F\subseteq G$, dan
$\mu F<\infty$, sehingga

\Centerline{$\lambda G\ge\nu F-\mu F
=\sum_{n=0}^m(\nu F_n-\mu F_n)\ge\gamma$.}

\noindent Karena $\gamma$ sembarang,
$\lambda G\ge\sum_{n=0}^{\infty}\lambda G_n$; karena
$\sequencen{G_n}$ sembarang, $\lambda$ aditif terhitung.\ \Qed

\medskip

\qquad\grheadb\ Sekarang $\mu+\lambda=\nu$. \Prf\ Domain dari
$\mu+\lambda$ adalah
$\dom\mu\cap\dom\lambda=\Tau=\dom\nu$. Ambil $G\in\Tau$. Jika
$\mu G=\infty$, maka $\nu G=\infty=(\mu+\lambda)G$. Jika tidak,

\Centerline{$(\mu+\lambda)G\ge\mu G+\nu G-\mu G=\nu G$.}

\noindent Jadi, jika $\nu G=\infty$, kita tentu akan mempunyai
$\nu G=(\mu+\lambda)G$. Terakhir, jika $\nu G<\infty$, maka

$$\eqalign{(\mu+\lambda)G
&=\mu G+\sup\{\nu F-\mu F:F\in\Tau,\,F\subseteq G\}\cr
&=\sup\{\nu F+\mu(G\setminus F):F\in\Tau,\,F\subseteq G\}\cr
&\le\sup\{\nu F+\nu(G\setminus F):F\in\Tau,\,F\subseteq G\}
=\nu G,\cr}$$

\noindent sehingga sekali lagi kita memperoleh kesamaan.\ \Qed

Dengan demikian, kita memperoleh penguraian $\nu$ yang diinginkan sebagai
suatu jumlah ukuran.

\medskip

{\bf (c)(i)} Jika $f$ nonnegatif, gabungkan (b) dan 234Hc.

\medskip

\quad{\bf (ii)} Secara umum, jika $\int fd\nu$ terdefinisi, demikian pula
$\int f^+d\nu$ dan $\int f^-d\nu$, dan paling banyak salah satunya
tak berhingga; sehingga $\int f^+d\mu$ dan $\int f^-d\mu$ terdefinisi
dan paling banyak salah satunya tak berhingga.
}%end of proof of 234Q

\exercises{
\leader{234X}{Latihan dasar (a)}
%\spheader 234Xa
\dvAformerly{1{}32Xk, 2{}14Xp}
Misalkan $(X,\Sigma,\mu)$ dan $(Y,\Tau,\nu)$ ruang-ruang ukur, dan
$\phi:X\to Y$ suatu fungsi \imp. Misalkan $A\subseteq X$ suatu himpunan
dengan ukuran luar penuh dalam $X$. Tunjukkan bahwa $\phi[A]$ mempunyai
ukuran luar penuh dalam $Y$, dan bahwa $\phi\restr A$ bersifat \imp\ 
untuk ukuran-ukuran subruang pada $A$ dan $\phi[A]$.
%234B

\spheader 234Xb\dvAformerly{1{}12Xg}
Misalkan $(X,\Sigma,\mu)$ suatu ruang ukur, $Y$ suatu himpunan, dan
$\phi:X\to Y$ suatu fungsi. Tunjukkan bahwa jika $\mu$ bertumpu pada
titik, maka demikian pula ukuran citra $\mu\phi^{-1}$.
%234D

\spheader 234Xc\dvAnew{2008}
Berikan suatu contoh ruang peluang $(X,\Sigma,\mu)$, suatu himpunan
$Y$, dan suatu fungsi $\phi:X\to Y$ sedemikian sehingga pelengkapan
ukuran citra $\mu\phi^{-1}$ bukan citra dari pelengkapan $\mu$.
\Hint{$\#(X)=3$.}
%234E

\spheader 234Xd\dvAnew{2008}
 Misalkan $X$, $Y$ himpunan-himpunan, $\phi:X\to Y$ suatu fungsi, dan
$\familyiI{\mu_i}$ suatu keluarga ukuran pada $X$ dengan jumlah $\mu$.
Dengan menuliskan $\mu_i\phi^{-1}$, $\mu\phi^{-1}$ untuk ukuran-ukuran
citra pada $Y$, tunjukkan bahwa
$\mu\phi^{-1}=\sum_{i\in I}\mu_i\phi^{-1}$.
%234G

\spheader 234Xe\dvAnew{2008}
Misalkan $X$ suatu himpunan.
(i) Tunjukkan bahwa jika $\familyiI{\mu_i}$ suatu keluarga terhitung dari
ukuran-ukuran $\sigma$-hingga pada $X$, dan
$\mu=\sum_{i\in I}\mu_i$ semihingga, maka $\mu$ $\sigma$-hingga.
(ii) Tunjukkan bahwa jika $\familyiI{\mu_i}$ suatu keluarga ukuran
atomik murni pada $X$, dan $\mu=\sum_{i\in I}\mu_i$ semihingga, maka
$\mu$ atomik murni.
(iii) Tunjukkan bahwa jika $\familyiI{\mu_i}$ sembarang keluarga ukuran
yang bertumpu pada titik pada $X$, maka $\sum_{i\in I}\mu_i$ bertumpu
pada titik.
%234G

\sqheader 234Xf\dvAnew{2008; lihat 1{}12Xe}
Misalkan $X$ suatu himpunan, dan tuliskan $\Mu$ untuk himpunan semua
ukuran pada $X$. Untuk $\mu\in\Mu$ dan
$\alpha\in\coint{0,\infty}$, definisikan $\alpha\mu$ dengan mengatakan
bahwa jika $\alpha>0$, maka
$(\alpha\mu)(E)=\alpha\mu E$ untuk $E\in\dom\mu$, sedangkan jika
$\alpha=0$, maka $(\alpha\mu)(E)=0$ untuk setiap $E\subseteq X$.
(i) Tunjukkan bahwa $\alpha\mu\in\Mu$ untuk semua
$\alpha\in\coint{0,\infty}$ dan $\mu\in\Mu$.
(ii) Tunjukkan bahwa $(\alpha+\beta)\mu=\alpha\mu+\beta\mu$,
$\alpha(\beta\mu)=(\alpha\beta)\mu$,
$\alpha(\mu+\nu)=\alpha\mu+\alpha\nu$ untuk semua $\alpha$,
$\beta\in\coint{0,\infty}$ dan $\mu$, $\nu\in\Mu$.
%234G

\spheader 234Xg\dvAformerly{2{}12Xj}
 Misalkan $X$ suatu himpunan, dan $\familyiI{\mu_i}$ suatu keluarga
ukuran lengkap pada $X$ dengan jumlah $\mu$.
Tunjukkan bahwa suatu fungsi bernilai $[-\infty,\infty]$ yang dinotasikan
dengan $f$ dan didefinisikan pada suatu subhimpunan dari $X$ terintegralkan terhadap $\mu$
jika dan hanya jika fungsi tersebut terintegralkan terhadap $\mu_i$ untuk
setiap $i\in I$ dan $\sum_{i\in I}\int|f|d\mu_i$ berhingga.
%234H

\spheader 234Xh\dvAformerly{2{}34Xa}
Misalkan $\mu$ ukuran Lebesgue pada $[0,1]$, dan tetapkan
$f(x)=\bover1x$ untuk $x>0$. Misalkan $\nu$ ukuran integral tak tentu
yang bersesuaian. Tunjukkan bahwa domain $\nu$ sama dengan domain $\mu$.
Tunjukkan bahwa untuk setiap $\delta\in\ocint{0,\bover12}$ terdapat
suatu himpunan terukur $E$ sedemikian sehingga
$\mu E=\delta$, tetapi $\nu E=\bover1{\delta}$.
%234K

\spheader 234Xi\dvAnew{2008}
 Misalkan $(X,\Sigma,\mu)$ suatu ruang ukur. (i) Tunjukkan bahwa jika
$\nu_1$ dan $\nu_2$ adalah ukuran integral tak tentu terhadap $\mu$,
maka demikian pula $\nu_1+\nu_2$. (ii) Tunjukkan bahwa jika
$\familyiI{\nu_i}$ suatu keluarga terhitung dari ukuran-ukuran integral
tak tentu terhadap $\mu$, dan $\nu=\sum_{i\in I}\nu_i$ semihingga, maka
$\nu$ adalah ukuran integral tak tentu terhadap $\mu$.
%234L

\spheader 234Xj\dvAformerly{2{}34Xb}
 Misalkan $(X,\Sigma,\mu)$ suatu ruang ukur, dan $\nu$ suatu ukuran
integral tak tentu terhadap $\mu$. Tunjukkan bahwa jika $\mu$ atomik
murni, maka demikian pula $\nu$.
%234N

\spheader 234Xk\dvAformerly{2{}34Xc}
 Misalkan $\mu$ suatu ukuran yang bertumpu pada titik. Tunjukkan bahwa
setiap ukuran integral tak tentu terhadap $\mu$ bertumpu pada titik.
%234N

\spheader 234Xl\dvAnew{2008}
 Misalkan $X$ suatu himpunan, dan $\Mu$ himpunan ukuran pada $X$, dengan
urutan parsial yang didefinisikan dalam 234P. Tunjukkan bahwa
(i) $\Mu$ mempunyai anggota terbesar dan terkecil (yang harus diuraikan);
(ii) jika $\familyiI{\mu_i}$ dan $\familyiI{\nu_i}$ adalah keluarga
dalam $\Mu$ sedemikian sehingga $\mu_i\le\nu_i$ untuk setiap $i$, maka
$\sum_{i\in I}\mu_i\le\sum_{i\in I}\nu_i$; (iii) jika kita mendefinisikan
perkalian skalar seperti dalam 234Xf, maka $\alpha\mu\le\mu$ kapan pun
$\mu\in\Mu$ dan $\alpha\in[0,1]$; (iv) dengan menuliskan $\hat\mu$ untuk
pelengkapan $\mu$, $\hat\mu\le\mu$ dan $\hat\mu\le\hat\nu$ kapan pun
$\mu$, $\nu\in\Mu$ dan $\mu\le\nu$; (v) dengan menuliskan $\tilde\mu$
untuk versi c.l.d.\ dari $\mu$, $\tilde\mu\le\mu$ untuk setiap
$\mu\in\Mu$; (vi) kapan pun $A\subseteq\Mu$ terarah ke atas, himpunan
tersebut mempunyai batas atas terkecil dalam $\Mu$.
%234Q

\spheader 234Xm\dvAnew{2008}
 Tuliskan pembuktian langsung elementer dari 234Qc yang tidak bergantung
pada 234Qb.
%234Q

\spheader 234Xn\dvAformerly{211Xd(v)} Misalkan $(X,\Sigma,\mu)$ dan
$(Y,\Tau,\nu)$ ruang-ruang ukur dan $\phi:X\to Y$ suatu fungsi \imp.
Tunjukkan bahwa jika $\mu$ $\sigma$-hingga dan atomik murni, maka
$\nu$ atomik murni.

\leader{234Y}{Latihan lanjutan (a)}
%\spheader 234Ya
\dvAformerly{2{}35Ya}
Tuliskan $\nu$ untuk ukuran Lebesgue pada $Y=[0,1]$, dan $\Tau$ untuk
domainnya. Misalkan $A\subseteq[0,1]$ suatu himpunan sedemikian sehingga
$\nu^*A=\nu^*([0,1]\setminus A)=1$, dan tetapkan
$X=[0,1]\cup\{x+1:x\in A\}\cup\{x+2:x\in[0,1]\setminus A\}$. Misalkan
$\mu_{LX}$ ukuran subruang yang diinduksi pada $X$ oleh ukuran Lebesgue,
dan tetapkan $\mu E=\bover13\mu_{LX}E$ untuk
$E\in\Sigma=\dom\mu_{LX}$. Definisikan $\phi:X\to Y$ dengan menuliskan
$\phi(x)=x$ jika $x\in[0,1]$, $\phi(x)=x-1$ jika
$x\in X\cap\ocint{1,2}$, dan $\phi(x)=x-2$ jika
$x\in X\cap\ocint{2,3}$. Tunjukkan bahwa $\nu$ adalah ukuran citra
$\mu\phi^{-1}$, tetapi $\nu^*A>\mu^*\phi^{-1}[A]$.
%234B

\spheader 234Yb\dvAnew{2008}
Carilah contoh-contoh menarik ruang peluang $(X,\Sigma,\mu)$ dan
$(Y,\Tau,\nu)$ yang mempunyai fungsi-fungsi $\phi:X\to Y$ sedemikian
sehingga $\phi[E]\in\Tau$ dan $\nu\phi[E]=\mu E$ untuk setiap
$E\in\Sigma$. \Hint{254K, 343J.}
%234B

\spheader 234Yc\dvAnew{2008}
Misalkan $\mu$ ukuran Lebesgue dua dimensi pada persegi satuan
$[0,1]^2$, dan misalkan $\phi:[0,1]^2\to[0,1]$ proyeksi pada koordinat
pertama, sehingga $\phi(\xi_1,\xi_2)=\xi_1$ untuk
$\xi_1$, $\xi_2\in[0,1]$. Tunjukkan bahwa ukuran citra
$\mu\phi^{-1}$ adalah ukuran Lebesgue satu dimensi pada $[0,1]$.
%234E

\spheader 234Yd\dvAformerly{1{}32Yf}
Dalam 234F, tunjukkan bahwa ukuran citra $\mu\phi^{-1}$ memperluas
$\nu$, dan sama dengan $\nu$ jika dan hanya jika $F\in\Tau$ untuk
setiap $F\subseteq Y\setminus\phi[X]$.
%234F

\spheader 234Ye\dvAnew{2008}
 Misalkan $(Y,\Tau,\nu)$ suatu ruang ukur lengkap, $X$ suatu himpunan,
dan $\phi:X\to Y$ suatu surjeksi. Tetapkan

\Centerline{$\Sigma=\{E:E\subseteq X$, $\phi[E]\in\Tau$,
$\nu(\phi[E]\cap\phi[X\setminus E])=0\}$,
\quad$\mu E=\nu\phi[E]$ untuk $E\in\Sigma$.}

\noindent Tunjukkan bahwa $\mu$ adalah pelengkapan dari ukuran yang
dibangun melalui proses 234F.
%234F

\spheader 234Yf\dvAnew{2008}
Misalkan $X$ suatu himpunan, dan $\Mu$ himpunan ukuran pada $X$.
Tunjukkan bahwa $\Mu$, dengan penjumlahan sebagaimana didefinisikan untuk
dua ukuran oleh rumus-rumus 234G, adalah semigrup komutatif dengan
identitas; uraikan identitas tersebut.
%234G

\spheader 234Yg\dvAnew{2008}
Berikan suatu contoh himpunan $X$, ukuran-ukuran peluang
$\mu_1$, $\mu_2$ pada $X$, dan suatu himpunan $A\subseteq X$ sedemikian
sehingga $A$ terabaikan terhadap $\mu_1$ maupun $\mu_2$, tetapi tidak
terabaikan terhadap $\mu$, dengan $\mu=\mu_1+\mu_2$.
%234G 

\spheader 234Yh\dvAformerly{2{}14Ya}
Dalam 214O, tunjukkan bahwa jika kita menetapkan
$\nu E=\sup_{I\in\Cal I}\mu^*(E\cap I)$ untuk setiap $E\in\Sigma$, maka
$\nu$ adalah ukuran, sedangkan $\mu=\nu+\lambda$.
%234H

\spheader 234Yi\dvAformerly{2{}34Yd}
 Misalkan $(X,\Sigma,\mu)$ suatu ruang ukur semihingga tanpa atom dan
$\nu$ suatu ukuran integral tak tentu terhadap $\mu$. Tunjukkan bahwa
pernyataan-pernyataan berikut ekuivalen: (i) untuk setiap $\epsilon>0$
terdapat suatu $\delta>0$ sedemikian sehingga $\nu E\le\epsilon$ kapan
pun $\mu E\le\delta$ (ii) $\nu$ mempunyai suatu turunan Radon-Nikod\'ym
yang dapat dinyatakan sebagai jumlah dari suatu fungsi terbatas dan suatu
fungsi terintegralkan.
%234K

\spheader 234Yj\dvAformerly{2{}34Yf}
 Misalkan $(X,\Sigma,\mu)$ suatu ruang ukur dan $\nu$ suatu ukuran
integral tak tentu terhadap $\mu$, dengan turunan Radon-Nikod\'ym $f$.
Tunjukkan bahwa versi c.l.d.\ dari $\nu$ adalah ukuran integral tak tentu
yang didefinisikan oleh $f$ terhadap versi c.l.d.\ dari $\mu$.
%234N

\spheader 234Yk\dvAformerly{2{}34Ya}
Misalkan $(X,\Sigma,\mu)$ suatu ruang ukur semihingga yang tidak dapat
dilokalkan. Tunjukkan bahwa terdapat suatu ukuran
$\nu:\Sigma\to[0,\infty]$ sedemikian sehingga $\nu E\le\mu E$ untuk
setiap $E\in\Sigma$, tetapi tidak terdapat fungsi terukur $f$ sedemikian
sehingga $\nu E=\int_Efd\mu$ untuk setiap $E\in\Sigma$.
%234O

\spheader 234Yl\dvAformerly{2{}34Yb}
Misalkan $(X,\Sigma,\mu)$ suatu ruang ukur yang dapat dilokalkan dengan
himpunan-himpunan terabaikan yang ditentukan secara lokal. Tunjukkan
bahwa suatu ukuran $\nu$, dengan domain $\Tau\supseteq\Sigma$, adalah
ukuran integral tak tentu terhadap $\mu$ jika dan hanya jika
($\alpha$) $\nu$ lengkap dan semihingga serta nol pada
himpunan-himpunan yang terabaikan terhadap $\mu$
($\beta$) kapan pun $\nu E>0$, terdapat suatu $F\subseteq E$ sedemikian
sehingga $F\in\Sigma$, $\mu F<\infty$, dan $\nu F>0$.
%234O

\spheader 234Ym\dvAformerly{2{}34Yc}
 Berikan suatu contoh ruang ukur yang dapat dilokalkan
$(X,\Sigma,\mu)$ dan suatu ukuran lengkap semihingga $\nu$ pada $X$,
yang didefinisikan pada suatu aljabar-sigma
$\Tau\supseteq\Sigma$, nol pada himpunan-himpunan yang terabaikan terhadap
$\mu$, dan sedemikian sehingga kapan pun $\nu E>0$, terdapat suatu
$F\subseteq E$ sedemikian sehingga $F\in\Sigma$, $\mu F<\infty$, dan
$\nu F>0$, tetapi $\nu$ bukan ukuran integral tak tentu terhadap $\mu$.
\Hint{216Yb.}
%234O 23bits

\spheader 234Yn\dvAformerly{2{}34Ye}
 Misalkan $(X,\Sigma,\mu)$ suatu ruang ukur yang dapat dilokalkan, dan
$\nu$ suatu ukuran lengkap yang dapat dilokalkan pada $X$, dengan domain
$\Tau\supseteq\Sigma$, yang merupakan pelengkapan dari pembatasannya pada
$\Sigma$. Tunjukkan bahwa jika kita menetapkan
$\nu_1F=\sup\{\nu(F\cap E):E\in\Sigma$, $\mu E<\infty\}$ untuk setiap
$F\in\Tau$, maka $\nu_1$ adalah ukuran integral tak tentu terhadap $\mu$,
dan terdapat suatu $H\in\Sigma$ sedemikian sehingga
$\nu_1F=\nu(F\cap H)$ untuk setiap $F\in\Tau$.
%234O

\spheader 234Yo\dvAnew{2008}
Misalkan $X$ suatu himpunan, dan $\Mu_{\text{sf}}$ himpunan ukuran
semihingga pada $X$. Untuk $\mu$, $\nu\in\Mu_{\text{sf}}$, katakan bahwa
$\mu\preccurlyeq\nu$ jika $\dom\nu\subseteq\dom\mu$,
$\mu F\le\nu F$ untuk setiap $F\in\dom\nu$, dan kapan pun
$E\in\dom\mu$ dan $\mu E>0$, terdapat suatu $F\in\dom\nu$ sedemikian
sehingga $F\subseteq E$ dan $0<\mu F<\infty$. (i) Tunjukkan bahwa
$(\Mu_{\text{sf}},\preccurlyeq)$ adalah suatu himpunan terurut parsial.
(ii) Tunjukkan bahwa jika $A\subseteq\Mu_{\text{sf}}$ suatu himpunan
tak kosong dengan batas atas dalam $\Mu_{\text{sf}}$, maka himpunan itu
mempunyai batas atas terkecil $\lambda$ yang didefinisikan dengan
mengatakan bahwa $\dom\lambda=\bigcap_{\mu\in A}\dom\mu$ dan, untuk
$E\in\dom\lambda$,

$$\eqalign{\lambda E
&=\sup\{\sum_{i=0}^n\mu_iF_i:\mu_0,\ldots,\mu_n\in A,\,
   \langle F_i\rangle_{i\le n}\text{ adalah partisi dari }E,\cr
&\mskip200mu
   F_i\in\dom\lambda\text{ untuk setiap }i\le n\}\cr
&=\sup\{\sum_{i=0}^n\mu_iF_i:\mu_0,\ldots,\mu_n\in A,\,
   F_0,\ldots,F_n\text{ saling lepas},\cr
&\mskip200mu
   F_i\in\dom\mu_i\text{ dan }F_i\subseteq E\text{ untuk setiap }i\le n\}.
   \cr}$$

\noindent(iii) Andaikan $\mu$, $\nu\in\Mu_{\text{sf}}$ mempunyai
pelengkapan $\hat\mu$, $\hat\nu$ dan versi-versi c.l.d.\ 
$\tilde\mu$, $\tilde\nu$.
Tunjukkan bahwa $\tilde\mu\preccurlyeq\hat\mu\preccurlyeq\mu$.
Tunjukkan bahwa jika $\mu\preccurlyeq\nu$, maka
$\hat\mu\preccurlyeq\hat\nu$ dan
$\tilde\mu\preccurlyeq\tilde\nu$.
%234P
}%end of exercises

\endnotes{\Notesheader{234}
Salah satu ciri mencolok teori ukuran, dibandingkan dengan cabang-cabang
matematika murni lain yang tingkat keabstrakannya sebanding, ialah relatif
tidak pentingnya konsep `morfisme' apa pun. Teori grup, misalnya,
didominasi oleh konsep `homomorfisme', dan topologi umum memberikan
tempat serupa kepada `fungsi kontinu'. Menurut saya, padanan yang
paling dekat dalam teori ukuran adalah gagasan `fungsi \imp' (234A).
Dalam Jilid 3 dan 4 saya bermaksud mengkaji konsep ini secara lebih
mendalam. Dalam jilid ini saya akan berpuas diri dengan menandai
fungsi-fungsi semacam itu ketika muncul dan dengan fakta-fakta dasar yang
tercantum dalam 234B.

Konstruksi `ukuran citra' (234C) secara alami terkait dengan
gagasan fungsi \imp. Ukuran semacam ini muncul di mana-mana dalam bidang
ini, dimulai dengan 234Yc yang tidak sepenuhnya elementer. Ukuran-ukuran
ini begitu penting sehingga wajar untuk mengkaji variasi-variasinya, seperti
dalam 234F dan 234Yb, tetapi menurut saya tidak satu pun memiliki arti
penting yang sebanding.

Hampir separuh bagian ini digunakan untuk `ukuran integral tak
tentu'. Saya menangani bagian ini dengan sangat hati-hati karena
gagasan-gagasan yang ingin saya nyatakan di sini, sejauh memperluas
pekerjaan dalam \S232, sangat bergantung pada perincian perumusan dalam
234I, dan mudah sekali mengambil langkah keliru setelah kita meninggalkan
konteks ukuran lengkap $\sigma$-hingga yang relatif terlindung. Saya
percaya bahwa jika kita bersedia mencurahkan sedikit usaha pada titik ini,
kita dapat mengembangkan suatu teori
(234K-234N) %234K 234L 234M 234N
yang akan menyediakan jalan mulus menuju penerapan-penerapan berikutnya;
untuk melihat maksud saya, Anda dapat merujuk pada entri-entri di bawah
`ukuran integral tak tentu' dalam indeks. Untuk saat ini saya hanya
menyebutkan suatu jenis teorema Radon-Nikod\'ym untuk ukuran-ukuran yang
dapat dilokalkan (234O).

Urutan parsial yang diuraikan dalam 234P-234Q hanyalah salah satu dari
banyak urutan yang dapat dipertimbangkan, dan untuk beberapa tujuan
tampaknya tidak memuaskan. Contoh-contoh terpenting akan muncul dalam
Bab 41 Jilid 4, dan mempunyai beragam ciri khusus yang mungkin patut
digunakan sebagai dasar bagi abstraksi lebih lanjut. Namun, versi di sini
memiliki keunggulan berupa kesederhanaan dan mendukung setidak-tidaknya
sebagian gagasan yang relevan (234Xl). Untuk konsep alternatif, lihat
234Yo.

}%end of notes

%look for ideas to add to \S412

\discrpage
